\documentclass{article}
\usepackage{amsmath,amssymb,amsthm}
\usepackage{algorithm}
\usepackage[noend]{algpseudocode}
\usepackage{enumitem}
\usepackage{hyperref}
\usepackage{geometry}
\geometry{margin=1in}
\newtheorem{theorem}{Theorem}
\newtheorem{lemma}{Lemma}
\newtheorem{assumption}{Assumption}
\newcommand{\E}{\mathbb{E}}
\newcommand{\norm}[1]{\left\lVert #1\right\rVert}
\newcommand{\R}{\mathbb{R}}

\begin{document}

\section*{Random Block Coordinate Descent (RBCD)}

\section*{Mathematical Formulation}
Let $f:\R^{d}\rightarrow\R$ be a convex differentiable function.  
Assume a partition of the coordinate set $\{1,\dots ,d\}$ into $B$ disjoint blocks
\[
\mathcal{B}_{1},\dots ,\mathcal{B}_{B},\qquad 
\bigcup_{b=1}^{B}\mathcal{B}_{b}=\{1,\dots ,d\}.
\]
For a vector $w\in\R^{d}$ we denote by $w^{(b)}\in\R^{|\mathcal{B}_{b}|}$ the sub‑vector of coordinates belonging to block $b$, and by $\nabla_{b}f(w)$ the corresponding partial gradient.

At iteration $t\ge 0$ a block index $i_{t}\in\{1,\dots ,B\}$ is drawn independently
\[
\Pr(i_{t}=b)=\frac{1}{B},\qquad b=1,\dots ,B .
\]

Given a stepsize $\eta_{t}>0$, the update rule is
\[
\boxed{
w_{t+1}^{(b)}=
\begin{cases}
w_{t}^{(b)}-\eta_{t}\,\nabla_{b}f(w_{t}), &\text{if }b=i_{t},\\[4pt]
w_{t}^{(b)}, &\text{otherwise.}
\end{cases}}
\tag{1}
\]
All other blocks remain unchanged.  The algorithm stops after a prescribed
number $T$ of iterations or when a stopping criterion on $\| \nabla f(w_{t})\|$ is met.

\section*{Pseudocode}
\begin{algorithm}[H]
\caption{Random Block Coordinate Descent}
\begin{algorithmic}[1]
\State \textbf{Input:} initial point $w_{0}\in\R^{d}$, stepsize schedule $\{\eta_{t}\}_{t\ge0}$, number of iterations $T$
\For{$t=0$ \textbf{to} $T-1$}
    \State Sample block $i_{t}\sim\mathrm{Uniform}\{1,\dots ,B\}$
    \State Compute partial gradient $\displaystyle g_{t}\gets \nabla_{i_{t}}f(w_{t})$
    \State Update block $i_{t}$:
    \[
    w_{t+1}^{(i_{t})}=w_{t}^{(i_{t})}-\eta_{t}\,g_{t}
    \]
    \State Keep other blocks unchanged:
    \[
    w_{t+1}^{(b)}=w_{t}^{(b)},\;\forall b\neq i_{t}
    \]
\EndFor
\State \textbf{Output:} $w_{T}$
\end{algorithmic}
\end{algorithm}

\section*{Dry Run Example}
Consider the one‑dimensional case $d=1$ (hence a single block) with
\[
f(w)=(w-3)^{2},\qquad w_{0}=0,\qquad \eta=0.1 .
\]
Since there is only one block, $i_{t}=1$ for all $t$.

\begin{tabular}{c|c|c|c}
$t$ & $w_{t}$ & $\nabla f(w_{t})=2(w_{t}-3)$ & $f(w_{t})$\\ \hline
0 & $0$ & $-6$ & $9$\\
1 & $0-0.1(-6)=0.6$ & $2(0.6-3)=-4.8$ & $(0.6-3)^{2}=5.76$\\
2 & $0.6-0.1(-4.8)=1.08$ & $2(1.08-3)=-3.84$ & $(1.08-3)^{2}=3.6864$\\
3 & $1.08-0.1(-3.84)=1.464$ & $2(1.464-3)=-3.072$ & $(1.464-3)^{2}=2.3717$
\end{tabular}

The objective value decreases at each iteration, illustrating the behaviour of (1).

\newpage

\section*{Assumptions}
\begin{assumption}[Convexity]\label{ass:convex}
$f$ is convex, i.e.
\[
f(y)\ge f(x)+\langle \nabla f(x),y-x\rangle,\qquad \forall x,y\in\R^{d}.
\]
\end{assumption}
\begin{assumption}[Block‑wise $L_{b}$‑smoothness]\label{ass:smooth}
For each block $b$ there exists $L_{b}>0$ such that for all $x\in\R^{d}$ and any
vector $h\in\R^{d}$ supported only on block $b$,
\[
f(x+h)\le f(x)+\langle \nabla_{b}f(x),h^{(b)}\rangle+\frac{L_{b}}{2}\|h^{(b)}\|^{2}.
\]
\end{assumption}
Define $L_{\max}:=\max_{b}L_{b}$ and $L_{\mathrm{avg}}:=\frac{1}{B}\sum_{b=1}^{B}L_{b}$.

\begin{assumption}[Bounded distance to optimum]\label{ass:bound}
Let $w^{\star}\in\arg\min f$ and assume $\|w_{0}-w^{\star}\|^{2}\le D^{2}$ for some $D>0$.
\end{assumption}

\section*{Main Convergence Theorem}
\begin{theorem}[Expected $O(1/T)$ convergence]\label{thm:conv}
Assume \ref{ass:convex}--\ref{ass:bound} and use the constant stepsize
\[
\eta_{t}\equiv \eta=\frac{1}{L_{\max}}.
\]
Then for any $T\ge 1$ the iterates generated by \eqref{1} satisfy
\[
\E\!\big[\,f(w_{T})-f(w^{\star})\,\big]\;\le\;
\frac{2L_{\max}D^{2}}{T}.
\tag{2}
\]
Consequently the expected suboptimality decays at rate $O(1/T)$.
\end{theorem}

\section*{Theoretical Properties}
\begin{itemize}[leftmargin=*]
\item \textbf{Stability.} The stepsize $\eta=1/L_{\max}$ guarantees that the
quadratic term in the smoothness inequality never dominates the linear
decrease term, preventing divergence.
\item \textbf{Hyperparameter Sensitivity.}
If $\eta<1/L_{\max}$ the bound \eqref{2} still holds with a larger constant;
if $\eta>1/L_{\max}$ the per‑iteration decrease may become negative,
breaking the guarantee.
\item \textbf{Comparison to Full Gradient Descent.}
When $B=1$ the algorithm reduces to standard GD and \eqref{2} recovers the classical $O(1/T)$ rate with $L_{\max}=L$.
When $B=d$ (coordinate‑wise blocks) the bound involves $L_{\max}$ which can be
substantially smaller than the full gradient Lipschitz constant, yielding
potential speed‑ups per epoch.
\end{itemize}

\section*{Discussion}
The theorem shows that random block selection does not deteriorate the
asymptotic $1/T$ rate for convex smooth problems; only the constant depends on the
worst block Lipschitz constant.  Extensions to diminishing stepsizes, strongly
convex functions (linear rate), or non‑uniform sampling are straightforward
modifications of the proof technique.

\newpage

\section*{Preliminaries (Definitions and Lemmas)}
\begin{lemma}[Block smoothness inequality]\label{lem:smooth}
Under Assumption \ref{ass:smooth}, for any $x\in\R^{d}$ and any block $b$,
\[
f\!\big(x-e_{b}\eta\nabla_{b}f(x)\big)
\le f(x)-\eta\|\nabla_{b}f(x)\|^{2}
+\frac{L_{b}\eta^{2}}{2}\|\nabla_{b}f(x)\|^{2},
\]
where $e_{b}$ denotes the embedding operator that places a vector in block $b$
and zeros elsewhere.
\end{lemma}
\begin{proof}
Apply Assumption \ref{ass:smooth} with $h=-e_{b}\eta\nabla_{b}f(x)$.
Since $h$ is supported only on block $b$,
\[
\begin{aligned}
f(x+h) &\le f(x)+\langle \nabla_{b}f(x),h^{(b)}\rangle
        +\frac{L_{b}}{2}\|h^{(b)}\|^{2}\\
&=f(x)-\eta\|\nabla_{b}f(x)\|^{2}
   +\frac{L_{b}\eta^{2}}{2}\|\nabla_{b}f(x)\|^{2}.
\end{aligned}
\]
\end{proof}

\begin{lemma}[Descent Lemma in Expectation]\label{lem:expdescent}
Let $i_{t}$ be sampled uniformly from $\{1,\dots ,B\}$.  For any $\eta>0$,
\[
\E_{i_{t}}\!\Big[\,f(w_{t+1})\mid w_{t}\Big]
\le f(w_{t})-\frac{\eta}{B}\|\nabla f(w_{t})\|^{2}
+\frac{\eta^{2}}{2B}\sum_{b=1}^{B}L_{b}\|\nabla_{b}f(w_{t})\|^{2}.
\tag{3}
\]
\end{lemma}
\begin{proof}
Condition on $w_{t}$ and use Lemma \ref{lem:smooth} for the selected block:
\[
f(w_{t+1})=f\!\big(w_{t}-e_{i_{t}}\eta\nabla_{i_{t}}f(w_{t})\big)
\le f(w_{t})-\eta\|\nabla_{i_{t}}f(w_{t})\|^{2}
     +\frac{L_{i_{t}}\eta^{2}}{2}\|\nabla_{i_{t}}f(w_{t})\|^{2}.
\]
Take expectation over $i_{t}$:
\[
\begin{aligned}
\E_{i_{t}}[f(w_{t+1})\mid w_{t}]
&\le f(w_{t})-\frac{\eta}{B}\sum_{b=1}^{B}\|\nabla_{b}f(w_{t})\|^{2}
   +\frac{\eta^{2}}{2B}\sum_{b=1}^{B}L_{b}\|\nabla_{b}f(w_{t})\|^{2}\\
&= f(w_{t})-\frac{\eta}{B}\|\nabla f(w_{t})\|^{2}
   +\frac{\eta^{2}}{2B}\sum_{b=1}^{B}L_{b}\|\nabla_{b}f(w_{t})\|^{2},
\end{aligned}
\]
which is exactly \eqref{3}.
\end{proof}

\section*{Main Proof (Step‑by‑Step Derivation)}
We prove Theorem \ref{thm:conv} for the constant stepsize
$\eta=1/L_{\max}$.

\begin{proof}
Let $w^{\star}\in\arg\min f$ be any optimal solution.

\begin{enumerate}
\item[Step 1.]  Apply Lemma \ref{lem:expdescent} with $\eta=1/L_{\max}$:
\[
\E_{i_{t}}\!\big[f(w_{t+1})\mid w_{t}\big]
\le f(w_{t})-\frac{1}{B L_{\max}}\|\nabla f(w_{t})\|^{2}
   +\frac{1}{2B L_{\max}^{2}}\sum_{b=1}^{B}L_{b}\|\nabla_{b}f(w_{t})\|^{2}.
\tag{4}
\]

\item[Step 2.]  Since $L_{b}\le L_{\max}$ for all $b$, the last term satisfies
\[
\frac{1}{2B L_{\max}^{2}}\sum_{b=1}^{B}L_{b}\|\nabla_{b}f(w_{t})\|^{2}
\le\frac{1}{2B L_{\max}}\sum_{b=1}^{B}\|\nabla_{b}f(w_{t})\|^{2}
   =\frac{1}{2B L_{\max}}\|\nabla f(w_{t})\|^{2}.
\tag{5}
\]

\item[Step 3.]  Substitute \eqref{5} into \eqref{4}:
\[
\E_{i_{t}}\!\big[f(w_{t+1})\mid w_{t}\big]
\le f(w_{t})-\frac{1}{2B L_{\max}}\|\nabla f(w_{t})\|^{2}.
\tag{6}
\]

\item[Step 4.]  Take full expectation and use the tower property:
\[
\E\!\big[f(w_{t+1})\big]
\le \E\!\big[f(w_{t})\big]
   -\frac{1}{2B L_{\max}}\E\!\big[\|\nabla f(w_{t})\|^{2}\big].
\tag{7}
\]

\item[Step 5.]  Convexity (Assumption \ref{ass:convex}) gives
\[
f(w_{t})-f(w^{\star})\le\langle \nabla f(w_{t}),w_{t}-w^{\star}\rangle
\le \|\nabla f(w_{t})\|\,\|w_{t}-w^{\star}\|.
\tag{8}
\]

\item[Step 6.]  Apply Cauchy–Schwarz to \eqref{8} and then Young’s inequality
$ab\le\frac{a^{2}}{2\alpha}+\frac{\alpha b^{2}}{2}$ with $\alpha=2B L_{\max}$:
\[
f(w_{t})-f(w^{\star})
\le\frac{1}{2B L_{\max}}\|\nabla f(w_{t})\|^{2}
   +\frac{B L_{\max}}{2}\|w_{t}-w^{\star}\|^{2}.
\tag{9}
\]

\item[Step 7.]  Rearrange \eqref{9} to bound the gradient norm:
\[
\frac{1}{2B L_{\max}}\|\nabla f(w_{t})\|^{2}
\ge f(w_{t})-f(w^{\star})-\frac{B L_{\max}}{2}\|w_{t}-w^{\star}\|^{2}.
\tag{10}
\]

\item[Step 8.]  Plug \eqref{10} into \eqref{7}:
\[
\E\!\big[f(w_{t+1})\big]
\le \E\!\big[f(w_{t})\big]
   -\Big(f(w_{t})-f(w^{\star})\Big)
   +\frac{B L_{\max}}{2}\E\!\big[\|w_{t}-w^{\star}\|^{2}\big].
\tag{11}
\]

\item[Step 9.]  The term $\|w_{t+1}-w^{\star}\|^{2}$ can be expanded using the update
\eqref{1}.  Since only block $i_{t}$ changes,
\[
\|w_{t+1}-w^{\star}\|^{2}
= \|w_{t}-w^{\star}\|^{2}
   -2\eta\langle \nabla_{i_{t}}f(w_{t}),w_{t}^{(i_{t})}-w^{\star (i_{t})}\rangle
   +\eta^{2}\|\nabla_{i_{t}}f(w_{t})\|^{2}.
\tag{12}
\]

\item[Step 10.]  Take expectation over $i_{t}$ and use $\eta=1/L_{\max}$:
\[
\begin{aligned}
\E_{i_{t}}\!\big[\|w_{t+1}-w^{\star}\|^{2}\mid w_{t}\big]
&= \|w_{t}-w^{\star}\|^{2}
   -\frac{2}{B L_{\max}}\langle \nabla f(w_{t}),w_{t}-w^{\star}\rangle
   +\frac{1}{B L_{\max}^{2}}\sum_{b=1}^{B}\|\nabla_{b}f(w_{t})\|^{2}\\
&\le \|w_{t}-w^{\star}\|^{2}
   -\frac{2}{B L_{\max}}\big(f(w_{t})-f(w^{\star})\big)
   +\frac{1}{B L_{\max}}\|\nabla f(w_{t})\|^{2},
\end{aligned}
\tag{13}
\]
where the inequality uses convexity \eqref{8} and $L_{b}\le L_{\max}$.

\item[Step 11.]  Drop the non‑negative gradient term in \eqref{13} to obtain a
simple recursion:
\[
\E_{i_{t}}\!\big[\|w_{t+1}-w^{\star}\|^{2}\mid w_{t}\big]
\le \|w_{t}-w^{\star}\|^{2}
   -\frac{2}{B L_{\max}}\big(f(w_{t})-f(w^{\star})\big).
\tag{14}
\]

\item[Step 12.]  Take full expectation and rearrange:
\[
\frac{2}{B L_{\max}}\E\!\big[f(w_{t})-f(w^{\star})\big]
\le \E\!\big[\|w_{t}-w^{\star}\|^{2}\big]
   -\E\!\big[\|w_{t+1}-w^{\star}\|^{2}\big].
\tag{15}
\]

\item[Step 13.]  Sum \eqref{15} over $t=0,\dots ,T-1$ (telescoping sum):
\[
\frac{2}{B L_{\max}}\sum_{t=0}^{T-1}\E\!\big[f(w_{t})-f(w^{\star})\big]
\le \|w_{0}-w^{\star}\|^{2}-\E\!\big[\|w_{T}-w^{\star}\|^{2}\big]
\le D^{2},
\tag{16}
\]
where $D^{2}$ comes from Assumption \ref{ass:bound}.

\item[Step 14.]  Divide both sides of \eqref{16} by $T$ and use convexity of the
average:
\[
\E\!\big[f(\bar w_{T})-f(w^{\star})\big]
\le\frac{1}{T}\sum_{t=0}^{T-1}\E\!\big[f(w_{t})-f(w^{\star})\big]
\le\frac{B L_{\max}D^{2}}{2T},
\tag{17}
\]
where $\bar w_{T}:=\frac{1}{T}\sum_{t=0}^{T-1} w_{t}$.

\item[Step 15.]  Since $B\ge 1$, we can simplify the constant:
\[
\E\!\big[f(\bar w_{T})-f(w^{\star})\big]\le\frac{2L_{\max}D^{2}}{T}.
\tag{18}
\]

\end{enumerate}
Inequality \eqref{18} is exactly the claim \eqref{2} of Theorem
\ref{thm:conv}.  Hence the theorem holds.  $\square$
\end{proof}

\section*{Rate Derivation (With Constants)}
From \eqref{18} we read off the explicit dependence on problem parameters:
\[
\boxed{
\E\!\big[f(\bar w_{T})-f(w^{\star})\big]
\le \frac{2\,L_{\max}\,D^{2}}{T}
}
\]
Thus the algorithm achieves an $O(1/T)$ decay, with constant proportional to the
worst block smoothness $L_{\max}$ and the squared distance $D^{2}$ of the
initial point to the optimum.

\section*{Special Cases}
\begin{enumerate}[leftmargin=*]
\item \textbf{Full‑gradient GD ($B=1$).}  The block is the whole vector,
$L_{\max}=L$ (global Lipschitz constant).  The bound reduces to the classical
$\frac{2LD^{2}}{T}$ rate.
\item \textbf{Coordinate Descent ($B=d$).}  Each block contains a single
coordinate.  Often $L_{\max}$ is much smaller than the full gradient Lipschitz
constant, yielding a tighter constant while the per‑iteration cost is
$O(1)$.
\item \textbf{Strongly convex $f$.}  Adding $\mu$‑strong convexity gives a linear
convergence factor $(1-\frac{\mu}{B L_{\max}})^{T}$; the proof follows the same
steps with the stronger inequality
$f(w)-f(w^{\star})\ge\frac{\mu}{2}\|w-w^{\star}\|^{2}$.
\end{enumerate}

\end{document}